\section{Scope, Status, and Qualifications}

\subsection{Governing Confession and Audience}

This monograph is an affirmative Catholic biblical, patristic, liturgical, and dogmatic study written for serious Catholic lay readers, catechists, seminarians, clergy, and study groups. It receives Sacred Scripture as the inspired word of God, reads the Old and New Testaments within the unity of God's saving economy and living Tradition (\textit{Dei Verbum} 11--12), and receives the four Marian dogmas according to the authoritative teaching of the Catholic Church. Jesus Christ alone is the eternal Word, the New Covenant, and the one Mediator between God and mankind (1 Tim 2:5; \textit{Lumen gentium} 60--62; CCC 970). Mary is confessed as the living Ark of the New Covenant because she truly conceived and bore the Son of God in her flesh; all her grace, dignity, and maternal intercession remain created, received from Christ, and wholly ordered to him.

\subsection{Biblical Witnesses, Book Naming, and Versification}

English Scripture quotations throughout this monograph reproduce verbatim the public-domain \textbf{Douay-Rheims Bible (Challoner revision)} held in the repository's source library, checked against the 1899 American Douay-Rheims edition and the \textbf{Clementine Vulgate (1598)}. Where a passage's modern English numbering or title differs from the traditional Vulgate and Douay-Rheims naming, standard modern titles (1--2 Samuel, 1--2 Kings, 1--2 Chronicles, Joshua, Jeremiah, 2 Maccabees, Revelation) are used in running prose alongside explicit Douay-Rheims counterparts where first needed (1--2 Kings = 1--2 Samuel; 3 Kings = 1 Kings; 1--2 Paralipomenon = 1--2 Chronicles; Josue = Joshua; Apocalypse = Revelation). Psalm citations give both the Vulgate/LXX number and the Hebrew number---for example, Psalm 131(132):8---and Exodus 40 notes the two-verse Vulgate offset (40:32--36 = modern 40:34--38).

Textual and lexical controls rely strictly on witnesses verified in the repository source library as of 25 November 2026:
\begin{itemize}
  \item \textbf{Greek New Testament:} Faithlife SBLGNT v1.2 (for Luke and Revelation), Robinson--Pierpont Byzantine Textform RP2018 (for Matthew, Mark, Luke, Acts, and Revelation), and the complete 27-book MorphGNT-SBLGNT v6.12 lemma-field projection (137,554 tokens) for exact New Testament counts of \textit{episkiazo} (5), \textit{anaphoneo} (1), \textit{skirtao} (3), \textit{kibotos} (6), \textit{kataskiazo} (1), and \textit{charitoo} (2). No Greek surface file for Hebrews is locally tracked; lemmas at Hebrews 9:4--5 are verified from the MorphGNT projection.
  \item \textbf{Greek Old Testament (Septuagint):} H. B. Swete's Greek text (First1KGreek digitisation) is locally tracked for Psalms and 3--4 Kingdoms (1--2 Kings), verifying \textit{episkiasei soi} at Ps 90:4, \textit{skirtao} at Ps 113:4, 6, \textit{he kibotos tou hagiasmatos sou} at Ps 131:8, and 3 Kingdoms 6 and 8. For Exodus, Numbers, Judges, 1--2 Samuel, 1--2 Chronicles, Jeremiah, and 2 Maccabees, Lancelot Brenton's public-domain English Septuagint translation (1851) is the tracked local witness; full Greek LXX morphological archives remain access-restricted outside Git and are never claimed as locally replayed counts.
  \item \textbf{Masoretic-Base English Controls:} The Catholic World English Bible (WEB), Revised Version (1895), and King James Version (1769) control Masoretic divergences at Judges 20:18, 26 (Bethel vs. Vulgate's Shiloh gloss), 1 Samuel 5:10 (Ekron vs. Brenton LXX Ashkelon), 2 Samuel 6:2 (Baale of Judah vs. Vulgate \textit{de viris Juda}), and Psalm 132:8 (``strength'' vs. LXX/Vulgate ``sanctification'').
  \item \textbf{Liturgical Propers:} Verified directly from the repository's normalized 1962 \textit{Missale Romanum} and postconciliar \textit{Missale Romanum} / \textit{Ordo lectionum Missae} (1981) calendar records, alongside the 1862 Pustet chant collation, the Benziger revised New Roman Missal, and the 1910 \textit{Raccolta} for the Litany of Loreto. Both Roman liturgical recensions are documented neutrally by official title and date.
\end{itemize}

\subsection{Evidentiary Distinctions and Attribution Controls}

Every claim in the monograph respects the distinction among literal scriptural narrative, canonical typology, original-language lexical evidence, patristic and liturgical reception, dogmatic definition, non-definitive magisterial teaching, scholastic theology, and spiritual synthesis:
\begin{itemize}
  \item \textbf{Verbatim quotation discipline:} Direct block and inline quotations of Church Fathers, Doctors, and popes come exclusively from public-domain primary texts tracked in the repository source library (including ANF Vol. 1, NPNF I--II, Ambrose's \textit{In Lucam} CSEL/Migne texts, Bede's \textit{Explanatio Apocalypsis}, Bernard's \textit{Super Missus est}, Aquinas's \textit{Summa theologiae} and \textit{Catena Aurea}, Bellarmine on the Psalms, Cornelius a Lapide, Hefele's \textit{History of the Councils}, Trent's decrees, Pius IX's \textit{Ineffabilis Deus} [1854 Latin and 1855 Baltimore bilingual], and Leo XIII's \textit{Quamquam pluries} [1889 Latin]). Works that are cataloged or metadata-only in the source library because of copyright or archival bounds---including \textit{Munificentissimus Deus}, \textit{Dei Verbum}, \textit{Lumen gentium}, the \textit{Catechism of the Catholic Church}, \textit{Redemptoris Mater}, \textit{Redemptoris Custos}, \textit{Patris Corde}, the Dicastery for the Doctrine of the Faith's \textit{Mater Populi Fidelis} (approved by Pope Leo XIV on 7 October 2025; published 4 November 2025), and John of Damascus's \textit{Homilies on the Dormition}---are summarized accurately by exact document, date, and paragraph without manufactured verbatim quotation.
  \item \textbf{Patristic attribution honesty:} Spurious or disputed homilies frequently cited in popular literature---notably the Coptic encomium on the Papyrus of Turin attributed to St. Athanasius and the Annunciation homilies printed under St. Gregory Thaumaturgus---are not presented as authentic third- or fourth-century patristic witnesses.
  \item \textbf{Modern canonical synthesis distinguished from early exegesis:} Where contemporary Catholic biblical theology observes rich structural parallels between 2 Samuel 6 and Luke 1 (such as Mary's three-month stay beside Obed-edom's three months), the monograph notes plainly where early and early-modern commentators (such as St. Ambrose and Cornelius a Lapide) did not draw that specific link, and distinguishes shared Greek vocabulary (\textit{episkiazo}) from thematic parallels using distinct verbs (\textit{skirtao} vs. David's dancing).
  \item \textbf{Dogmatic and pastoral bounds:} Mary's cooperation in redemption is not proposed as a fifth Marian dogma, and Saint Joseph is never treated as a typological fulfillment of Uzzah.
\end{itemize}

\subsection{Artwork, Cartography, and Web-Edition Equivalence}

The six monochrome graphite illustrations in the print edition were produced specifically for this monograph under strict scriptural and visual controls documented in \texttt{research/artwork-audit.md}. All labels, callouts, station names, and route overlays are typeset live in TeX/TikZ over text-free monochrome PNG bases. The illustrations and topographic base map provide visual orientation and disciplined architectural imagination; they neither reconstruct lost ancient artifacts nor assert exact ancient roadbeds or archaeological proof for unlocated sites (such as the house of Obed-edom or the unnamed town of Judah in Luke 1:39). In the web edition, every graphical plate and TikZ overlay is replaced cleanly by its complete authored textual equivalent so that no substantive geographical, architectural, or theological content is lost in reflow.

\subsection{Use, Rights, and Review Status}

This publication is a study and devotional aid, not an official liturgical book, catechism, critical textual edition, or ecclesiastical judgment. It carries no diocesan \textit{imprimatur} and no external specialist or clerical review. Project-authored prose, design, and bespoke graphite illustrations are released under the repository license noted in the final colophon; incorporated public-domain Scripture, patristic translations, liturgical texts, and third-party lexical controls retain their own respective public-domain or Creative Commons statuses as recorded in \texttt{research/source-audit.md} and \texttt{research/source-bindings.toml}.
